\documentclass[11pt]{article}
\usepackage[a4paper,margin=18mm]{geometry}
\usepackage[no-math]{fontspec}
\setmainfont{Latin Modern Roman}
\usepackage{amsmath,amssymb,amsthm,xfrac,xspace,hyperref,graphicx,comment}
\excludecomment{DefTralics}

\providecommand{\bibliofont}{\small}
\newcommand{\ie}{i.e.,\xspace}
\newcommand{\eg}{e.g.,\xspace}
\newcommand{\etc}{etc.\xspace}
\newcommand{\cf}{cf.\xspace}
\newcommand{\resp}{resp.\xspace}
\newcommand{\smaller}{\small}
\newcommand{\Subsection}{\subsection}
\newcommand{\Subsubsection}{\subsubsection}
\newcommand{\skpt}{\smallskip}
\emergencystretch=3em
\numberwithin{equation}{section}\input{author-preamble.tex}
\begin{document}
\begin{center}{\Large Stable item 2041}\end{center}
\paragraph{Source and attribution.}
Tomoyuki Arakawa, Cuipo Jiang, Anne Moreau, \emph{Simplicity of vacuum modules and associated varieties}, Journal de l’École polytechnique — Mathématiques 8 (2021), publisher version, DOI 10.5802/jep.144. \href{https://jep.centre-mersenne.org/articles/10.5802/jep.144/}{Publisher article}. Authors' text: \href{https://creativecommons.org/licenses/by/4.0/}{CC BY 4.0}.
\paragraph{Editorial problem boundary.}
Prove only Theorem 1.2 below for a finite-dimensional complex simple Lie algebra and noncritical level \(k\ne-h^\vee\): every nonzero singular vector of the universal affine vertex algebra has nonzero image in its Zhu \(C_2\)-algebra. The entire authored commutative reduction and noncommutative Sugawara/PBW cancellation argument is retained, including its printed bounded ancillary PBW checks. Critical-level, superalgebra and arbitrary singular arc-polynomial assertions are outside this item.
\paragraph{Difficulty (editorial): H2.}
The nonvanishing of the Zhu image is not a formal consequence of the associated graded invariant ring, whose singular arc-polynomials can have positive depth. The author combines the commutative highest-weight reduction with a carefully chosen PBW monomial and an explicit Sugawara commutator expansion that controls its root and Cartan factors. The noncritical-level coefficient then rules out the proposed leading-term cancellation and forces zero depth, which is the nonroutine obligation behind the selected singular-vector statement.
\paragraph{Editorial transcription note.}
Original author language, mathematics and ordering are preserved. Publisher front matter is replaced by this independent article wrapper. Bibliography is recovered from matching cached publisher HTML, including its TeX formulas. No new proof or mathematical repair is supplied. Surrounding original results are retained for conservative context and are not extra selected corpus claims.
\clearpage
\input{author-body.tex}
\begin{thebibliography}{99}
\bibitem{AbeBuhDon04} T. Abe, G. Buhl \& C. Dong - “Rationality, regularity, and $C_2$-cofiniteness” , Trans. Amer. Math. Soc. 356 (2004) no. 8, p. 3391-3402
\bibitem{Arakawa-Kawasetsu} T. Arakawa \& K. Kawasetsu - “Quasi-lisse vertex algebras and modular linear differential equations” , in Lie groups, geometry, and representation theory , Progress in Math. , vol. 326 , Birkhäuser/Springer, Cham, 2018, p. 41-57
\bibitem{AM16} T. Arakawa \& A. Moreau - “Sheets and associated varieties of affine vertex algebras” , Adv. Math. 320 (2017), p. 157-209, Corrigendum: Ibid 372 (2020), article Id. 107302
\bibitem{AM15} T. Arakawa \& A. Moreau - “Joseph ideals and lisse minimal $W$-algebras” , J. Inst. Math. Jussieu 17 (2018) no. 2, p. 397-417
\bibitem{AM17} T. Arakawa \& A. Moreau - “On the irreducibility of associated varieties of $W$-algebras” , J. Algebra 500 (2018), p. 542-568
\bibitem{Ara05} T. Arakawa - “Representation theory of superconformal algebras and the Kac-Roan-Wakimoto conjecture” , Duke Math. J. 130 (2005) no. 3, p. 435-478
\bibitem{Ara07} T. Arakawa - “Representation theory of $\mathscr{W}$-algebras” , Invent. Math. 169 (2007) no. 2, p. 219-320
\bibitem{Ara08-a} T. Arakawa - “Representation theory of $W$-algebras, II” , in Exploring new structures and natural constructions in mathematical physics , Adv. Stud. Pure Math. , vol. 61 , Math. Soc. Japan, Tokyo, 2011, p. 51-90
\bibitem{Ara12} T. Arakawa - “A remark on the $C_2$-cofiniteness condition on vertex algebras” , Math. Z. 270 (2012) no. 1-2, p. 559-575
\bibitem{A11} T. Arakawa - “$W$-algebras at the critical level” , in Algebraic groups and quantum groups , Contemp. Math. , vol. 565 , American Mathematical Society, Providence, RI, 2012, p. 1-13
\bibitem{Ara09b} T. Arakawa - “Associated varieties of modules over Kac-Moody algebras and $C_2$-cofiniteness of $W$-algebras” , Internat. Math. Res. Notices (2015) no. 22, p. 11605-11666
\bibitem{A2012Dec} T. Arakawa - “Rationality of $W$-algebras: principal nilpotent cases” , Ann. of Math. (2) 182 (2015) no. 2, p. 565-604
\bibitem{AEkeren19} T. Arakawa \& J. van Ekeren - “Rationality and fusion rules of exceptional W-algebras” , 2019
\bibitem{BeiDri} A. Beilinson \& V. Drinfeld - “Quantization of Hitchin’s integrable system and Hecke eigensheaves” , preprint, available at http://math.uchicago.edu/\textasciitilde{}drinfeld/langlands/QuantizationHitchin.pdf
\bibitem{BeiFeiMaz} A. Beilinson, B. Feigin \& B. Mazur - “Introduction to algebraic field theory on curves” , preprint
\bibitem{BeeLemLie15} C. Beem, M. Lemos, P. Liendo, W. Peelaers, L. Rastelli \& B. C. van Rees - “Infinite chiral symmetry in four dimensions” , Comm. Math. Phys. 336 (2015) no. 3, p. 1359-1433
\bibitem{BeeRas} C. Beem \& L. Rastelli - “Vertex operator algebras, Higgs branches, and modular differential equations” , J. High Energy Phys. (2018) no. 8, article ID 114, 70 pages
\bibitem{Charbonnel-Moreau} J.-Y. Charbonnel \& A. Moreau - “The symmetric invariants of centralizers and Slodowy grading” , Math. Z. 282 (2016) no. 1-2, p. 273-339
\bibitem{DonMas06} C. Dong \& G. Mason - “Integrability of $C_2$-cofinite vertex operator algebras” , Internat. Math. Res. Notices (2006), article ID 80468, 15 pages
\bibitem{DSK06} A. De Sole \& V. G. Kac - “Finite vs affine $W$-algebras” , Japan. J. Math. 1 (2006) no. 1, p. 137-261
\bibitem{EisFre01} D. Eisenbud \& E. Frenkel - “Appendix to [Mus01]” , 2001
\bibitem{FeiFre90} B. Feigin \& E. Frenkel - “Quantization of the Drinfel’d-Sokolov reduction” , Phys. Lett. B 246 (1990) no. 1-2, p. 75-81
\bibitem{FeiFre92} B. Feigin \& E. Frenkel - “Affine Kac-Moody algebras at the critical level and Gel’fand-Dikiĭ algebras” , in Infinite analysis, Part A, B (Kyoto, 1991) , Adv. Ser. Math. Phys. , vol. 16 , World Sci. Publ., River Edge, NJ, 1992, p. 197-215
\bibitem{FreGai04} E. Frenkel \& D. Gaitsgory - “$D$-modules on the affine Grassmannian and representations of affine Kac-Moody algebras” , Duke Math. J. 125 (2004) no. 2, p. 279-327
\bibitem{Fre05} E. Frenkel - “Wakimoto modules, opers and the center at the critical level” , Adv. Math. 195 (2005) no. 2, p. 297-404
\bibitem{FZ} I. B. Frenkel \& Y. Zhu - “Vertex operator algebras associated to representations of affine and Virasoro algebras” , Duke Math. J. 66 (1992) no. 1, p. 123-168
\bibitem{GanGin02} W. L. Gan \& V. Ginzburg - “Quantization of Slodowy slices” , Internat. Math. Res. Notices (2002) no. 5, p. 243-255
\bibitem{GorKac07} M. Gorelik \& V. G. Kac - “On simplicity of vacuum modules” , Adv. Math. 211 (2007) no. 2, p. 621-677
\bibitem{Hartshorne} R. Hartshorne - Algebraic geometry , Graduate Texts in Math. , vol. 52 , Springer-Verlag, New York-Heidelberg, 1977
\bibitem{Hu72} J. E. Humphreys - Introduction to Lie algebras and representation theory , Graduate Texts in Math. , vol. 9 , Springer-Verlag, New York-Berlin, 1972
\bibitem{Kac_infinite} V. G. Kac - Infinite-dimensional Lie algebras , Cambridge University Press, Cambridge, 1990
\bibitem{KacRoaWak03} V. G. Kac, S.-S. Roan \& M. Wakimoto - “Quantum reduction for affine superalgebras” , Comm. Math. Phys. 241 (2003) no. 2-3, p. 307-342
\bibitem{KacWak89} V. G. Kac \& M. Wakimoto - “Classification of modular invariant representations of affine algebras” , in Infinite-dimensional Lie algebras and groups (Luminy-Marseille, 1988) , Adv. Ser. Math. Phys. , vol. 7 , World Sci. Publ., Teaneck, NJ, 1989, p. 138-177
\bibitem{KacWak08} V. G. Kac \& M. Wakimoto - “On rationality of $W$-algebras” , Transform. Groups 13 (2008) no. 3-4, p. 671-713
\bibitem{Li05} H. Li - “Abelianizing vertex algebras” , Comm. Math. Phys. 259 (2005) no. 2, p. 391-411
\bibitem{LL} J. Lepowsky \& H. Li - Introduction to vertex operator algebras and their representations , Progress in Math. , vol. 227 , Birkhäuser Boston, Inc., Boston, MA, 2004
\bibitem{Miy04} M. Miyamoto - “Modular invariance of vertex operator algebras satisfying $C_2$-cofiniteness” , Duke Math. J. 122 (2004) no. 1, p. 51-91
\bibitem{Mus01} M. Mustaţă - “Jet schemes of locally complete intersection canonical singularities” , Invent. Math. 145 (2001) no. 3, p. 397-424
\bibitem{Premet02} A. Premet - “Special transverse slices and their enveloping algebras” , Adv. Math. 170 (2002) no. 1, p. 1-55
\bibitem{RaiTau92} M. Raïs \& P. Tauvel - “Indice et polynômes invariants pour certaines algèbres de Lie” , J. reine angew. Math. 425 (1992), p. 123-140
\bibitem{Slo80} P. Slodowy - Simple singularities and simple algebraic groups , Lect. Notes in Math. , vol. 815 , Springer, Berlin, 1980
\bibitem{XieYan2019lisse} D. Xie \& W. Yan - “4d $\mathcal{N}=2$ SCFTs and lisse W-algebras” , 2019
\bibitem{Zhu96} Y. Zhu - “Modular invariance of characters of vertex operator algebras” , J. Amer. Math. Soc. 9 (1996) no. 1, p. 237-302
\end{thebibliography}
\end{document}
